% Chapter 1 --- Introduction

\chapter{Introduction}\doublespacing
\label{Chapter1}
\lhead{Chapter I. \emph{Introduction}}

%----------------------------------------------------------------------------------------
\section{Introduction}
%----------------------------------------------------------------------------------------

The first thing that strikes a visitor to an ITC procurement warehouse is not the scale of the operation---it is the silence of the sorting table. In the midst of arriving trucks, diesel engines, and shouted instructions, multiple technicians sit hunched over a 100-gram spread of wheat, working grain by grain with practiced hands. This act is the quality gate through which every consignment must pass before ITC's \textit{Aashirwad Atta} supply chain can proceed \cite{Reardon2009}.

Watching this process during initial fieldwork, several things became clear almost immediately. The task is extraordinarily difficult to standardise. The categories the assessor must distinguish---\textit{Potia}, \textit{Black Tip}, early \textit{Infested} grain---are not well-defined in written manuals; they live in the assessor's hands and eyes. And the time it takes, on average, 60 minutes per 100-gram sample, creates a bottleneck that no amount of additional personnel can fully resolve because the bottleneck is not headcount. It is the fundamental speed of careful human attention applied to 2,800 individual entities.

This thesis documents the development of the \textbf{``Drishti'' Wheat Grain Analyser}, an automated assessment device designed to address this constraint at its source. The name draws on the Sanskrit word for sight or vision---an appropriate choice given that the system's central challenge is not speed per se, but the quality of observation that speed must not compromise. The device was developed for and in dialogue with ITC Limited, whose procurement scale makes the latency problem both economically significant and technically demanding.

The research was based on the conviction that good design must begin with understanding before it begins with making. No prototype was sketched until substantial time had been spent in the field, in warehouses, at sorting tables, in conversations with assessors, quality managers, and the food analysts at ITC's Bangalore headquarters who ultimately define what ``quality'' means in this context. That orientation, field-first, shaped every technical and design decision that followed.

%----------------------------------------------------------------------------------------
\section{Problem Statement}
%----------------------------------------------------------------------------------------

A trained technician at an ITC procurement warehouse must sort a 100-gram wheat sample into 12 distinct categories: pure wheat, broken wheat, infested, shrivelled, blacktip, kernel burnt, potia, stones, chaff, mud balls, sticks, and other food grains. For each, the weight-by-weight percentage ($W/W\%$) is calculated and checked against acceptance thresholds before a truck is cleared for unloading. This is not a trivial categorisation task. At approximately 2,800 entities per sample, it requires close attention to morphological details, surface colour, kernel geometry, the faint darkening at a grain's tip that marks it as \textit{Black Tip}, that take months, sometimes years, of practice to read reliably.

The time implications are significant. At 60 minutes per cycle, assessment becomes the rate-limiting step in a continuously operating logistics chain. Trucks queue. Assessors, under pressure, tire. And the categories most sensitive to fatigue, the subtle ones like \textit{Potia} and early-stage \textit{Infested} grain, are precisely those where errors carry the most downstream consequence. The problem is not simply that the process is slow. It is that speed and accuracy are in structural tension, and the margin between them narrows as a shift progresses.

The problem, as distilled from extended fieldwork, is therefore compound: one of \textit{assessment latency}, \textit{quality inconsistency across operators and shifts}, and \textit{fragile concentration of tacit expertise} in a small number of senior assessors on whose presence the entire system depends.

%----------------------------------------------------------------------------------------
\section{Research Objectives and Scope}
%----------------------------------------------------------------------------------------

The overarching aim of this research is to design, build, and evaluate an automated wheat grain quality assessment system that is genuinely suited to the conditions of ITC's procurement environment, not an idealised version, but the real one, with its dust, its noise, and its time-pressured operators. From that aim, the following specific objectives were derived:

\begin{itemize}
    \item \textbf{Contextual inquiry of the procurement ecosystem:} To conduct structured field observation at both warehouse and factory sites, mapping the assessment workflow, documenting the record systems in active use, and surfacing the tacit knowledge that shapes how experienced assessors make category calls.

    \item \textbf{Derivation of functional and ergonomic constraints:} To establish, from field evidence rather than assumption, the physical and operational requirements that any viable automated system must satisfy: handling of dense 100-gram samples ($\sim$2,800 entities), dust tolerance, and usability under time pressure.

    \item \textbf{Design of a mechanical-computational intervention:} To develop a system architecture that addresses the fundamental occlusion challenge in dense grain imaging, not by asking the algorithm to compensate for poor observability, but by engineering the physical conditions of observation so that the algorithm's task becomes tractable.

    \item \textbf{Iterative prototype development and expert validation:} To build through successive functional prototypes, from small-scale proof-of-concept to field-ready unit, testing each generation against real grain samples and incorporating feedback from ITC's quality assurance team.

    \item \textbf{Contribution of design guidelines for industrial AI:} To document the design decisions made in this project in terms that are useful to practitioners working at the intersection of industrial design and applied machine learning.
\end{itemize}

The research is scoped to the procurement and storage phases of the ITC wheat supply chain, focusing on the twelve grain and impurity categories specified in ITC's quality standards. Processing and packaging stages fall outside its boundary.

%----------------------------------------------------------------------------------------
\section{Thesis Organisation}
%----------------------------------------------------------------------------------------

The thesis is structured to reflect how the research actually unfolded---from initial field immersion through to design synthesis and expert evaluation.

Chapter~\ref{Chapter2} presents the contextual inquiry: observations and interviews at ITC warehouses and factories that gave shape to the design brief. The reader who skips this chapter will not understand why subsequent design decisions were made.

Chapter~\ref{Chapter3} reviews the relevant literature - computer vision for grain classification, the specific failure modes of dense-object detection, and the hardware-software co-design framework that the research ultimately adopts.

Chapter~\ref{Chapter4} traces the design methodology and prototyping sequence, from early illumination experiments through two prototype generations to the pivot that defined the rest of the project.

Chapter~\ref{Chapter5} addresses the mechanical design problem in detail: why vibration failed, how electromagnetic solenoid actuation was developed as an alternative, and how the asynchronous pulse timing logic was arrived at empirically.

Chapter~\ref{Chapter6} covers the computational architecture: optoelectrical calibration, dataset construction strategy, and the edge-to-server deployment model.

Chapter~\ref{Chapter7} examines human-device interaction, addressing both the physical tray interface and the digital dashboard that communicates assessment outcomes to operators.

Chapter~\ref{Chapter8} reports expert validation outcomes from the demonstration conducted at ITC headquarters in Bangalore.

Chapter~\ref{Chapter9} synthesises the research contributions, acknowledges the study's limitations with appropriate candour, and identifies the most productive directions for future work.

%----------------------------------------------------------------------------------------
\section{Summary}
%----------------------------------------------------------------------------------------

The problem addressed in this thesis is deceptively simple to state: a quality assessment that takes sixty minutes needs to take five. The actual design challenge, as the fieldwork made clear, is considerably more complex. It involves understanding what experienced assessors actually do, and why, before proposing any alternative. The chapters that follow are an account of how that understanding was built and what it made possible.
